\documentclass[../main.tex]{subfiles}
\begin{document}

\chapter{Numerical experiments}
\label{chap:EXPERIMENTS}
This chapter consists of numerical experiments using the methods described in the previous chapters.\sidenote{The numerical experiments demonstrating the failure of SILVA to preserve hydrostatic balance are presented in \Cref{sec:SILVA_FAILURE} and not here, the reason being that we wish to present \textcolor{gray}{at least partially} successful results in this chapter.} Namely, we compare the classical SPH method \Cref{sec:CLASSIC-COMPLETE} and the pressure-potential temperature density-free formulation \Cref{sec:DI-pressure-theta} with their perturbation/well-balanced variants \Cref{sec:PERTURBATION-CLASSICAL-SPH} and \Cref{sec:WB-DI-PTH} respectively.

All numerical experiments presented in this work were conducted using the code described in \Cref{chap:IMPLEMENTATION}. Every simulation in this chapter is fully specified by its parameter files and can be regenerated from the archived code, see \Cref{sec:REPRODUCIBILITY}. The parameters are collected in \Cref{tab:PARAMS-PHYSICS,tab:PARAMS-GEOM,tab:PARAMS-SPH,tab:PARAMS-NUM}.

\begin{table}[htbp]
    \centering
    \caption{Physical constants.}
    \label{tab:PARAMS-PHYSICS}
    \begin{tblr}{
        colspec = {lc},
        row{1} = {guard, font=\bfseries\sffamily},
    }
        \toprule
        Parameter & Value \\
        \midrule
        Gravitational acceleration $g$ [\qty{}{\metre\per\second\squared}] & 9.81 \\
        Specific gas constant, dry air $R$ [\qty{}{\joule\per\kilogram\per\kelvin}] & 287.05 \\
        Adiabatic index $\gamma$ & 1.4 \\
        Reference temperature $T_{bg}$ [\qty{}{\kelvin}] & 250.0 \\
        Reference density $\rho_0$ [\qty{}{\kilogram\per\cubic\metre}] & 1.393 \\
        Brunt--V\"ais\"al\"a frequency $N$ [\qty{}{\per\second}] & 0.14 \\
        Initial (background) horizontal wind $U_0$ [\qty{}{\metre\per\second}] & 20.0 \\
        \midrule
        Scale height (derived) $H = R\,T_{bg}/g$ [\qty{}{\metre}] & 7315.2 \\
        \bottomrule
    \end{tblr}
\end{table}

\begin{table}[htbp]
    \centering
    \caption{Domain and mountain geometry.}
    \label{tab:PARAMS-GEOM}
    \begin{tblr}{
        colspec = {lc},
        row{1} = {guard, font=\bfseries\sffamily},
    }
        \toprule
        Parameter & Value \\
        \midrule
        Domain length $L_x$ [\qty{}{\metre}] & 400000.0 \\
        Domain height / top boundary $z_t$ [\qty{}{\metre}] & 26000.0 \\
        Witch of Agnesi half-width $a$ [\qty{}{\metre}] & 10000.0 \\
        Mountain height $h_m$ [\qty{}{\metre}] & 100.0 \\
        \bottomrule
    \end{tblr}
\end{table}

\begin{table}[htbp]
    \centering
    \caption{SPH discretization parameters, parameters of the quintic Wendland (2D).}
    \label{tab:PARAMS-SPH}
    \begin{tblr}{
        colspec = {lc},
        row{1} = {guard, font=\bfseries\sffamily},
    }
        \toprule
        Parameter & Value \\
        \midrule
        Particle spacing $\Delta r$ [\qty{}{\metre}] & 520.0 \\
        Smoothing length coefficient $\eta$ & 1.8 \\
        Artificial viscosity $\alpha$ & 0.15 \\
        Artificial viscosity $\beta$ & 0.3 \\
        XSPH velocity-smoothing coefficient $\epsilon$ & 0.01 \\
        \midrule
        Smoothing length (derived) $h_0 = \eta\,\Delta r$ [\qty{}{\metre}] & 936.0 \\
        \bottomrule
    \end{tblr}
\end{table}

\begin{table}[htbp]
    \centering
    \caption{Time-integration and damping parameters.}
    \label{tab:PARAMS-NUM}
    \begin{tblr}{
        colspec = {lcc},
        row{1} = {guard, font=\bfseries\sffamily},
    }
        \toprule
        Parameter & Stationary Case & Evolutionary Case \\
        \midrule
	Relative time step $\Delta t_{\text{rel}}$ & 0.1 & 0.1 \\
	End time $t_{\text{end}}$ [\qty{}{\second}] & 100.0 & 1200.0 \\
	Relative Rayleigh damping $\gamma_{r,\text{rel}}$ & 0.2 & 0.2 \\
        Damping layer thickness $z_\beta$ [\qty{}{\metre}] & 3000.0 & 3000.0 \\
        \midrule
        Reference sound speed $c_{ref}$ [\qty{}{\metre\per\second}] & 255.6 & 255.6 \\
        Time step $\Delta t$ [\qty{}{\second}] & 0.366 & 0.366 \\
	Number of time steps $t_{\text{end}}/\Delta t$ & $\approx 273$ & $\approx 3278$ \\
        Rayleigh damping rate $\gamma_r$ [\qty{}{\per\second}] & 0.028 & 0.0 \\
        Damping layer lower bound $z_t - z_\beta$ [\qty{}{\metre}] & 23000.0 & --- \\
        \bottomrule
    \end{tblr}
    \vspace{2pt}
{\footnotesize Derived quantities (lower half): $c_{\text{ref}} = \sqrt{\gamma P_{\text{ref}}/\rho_0}$ with fixed reference pressure $P_{\text{ref}} = \qty{65}{\kilo\pascal}$; $\Delta t = \Delta t_{\text{rel}}\,h_0/c_{\text{ref}}$; $\gamma_r = \gamma_{r,\text{rel}}\,N$; the damping layer occupies $z \in [z_t - z_\beta,\, z_t]$.}
\end{table}

\section{Preservation of hydrostatic balance (stationary case)}
\label{sec:STATIONARY}

As established in \Cref{sec:WB-WAVES}, preservation of the hydrostatic balance on \emph{the discrete level} is essential in simulations of internal gravity waves. In this section, we compare the ability of the used formulations to preserve the initial hydrostatic state.

To describe the velocity field quantitatively, we plot the \emph{maximum particle velocity magnitude}, that is

\begin{equations}[single,!,numberline=all]
\label{eq:L_infty}
\abs{\vb{v}}_{\infty} \coloneq \max_{\alpha\in \{1, \dots, N\}} \abs{\vb{v}_{\alpha}},
\end{equations}

and also the \emph{root-mean-square velocity} 

\begin{equations}[single,!,numberline=all]
\label{eq:L_2}
\abs{\vb{v}}_2 \coloneq \frac{1}{N}\sum_{\alpha=1}^N \abs{\vb{v}_{\alpha}}^{2}.
\end{equations}

The measures $\abs{\vb{v}}_{\infty}, \abs{\vb{v}}_{2}$ are discrete counterparts of the $L_{\infty}$ and $L_{2}$ norms.

Figures \Cref{fig:STATIONARY-NORMS} depict the time evolution of the velocity field of an atmosphere initialized from hydrostatic balance. They show the \emph{failure} of the standard formulations to \emph{preserve the hydrostatic balance}. On the contrary, the perturbation/well-balanced variants retain the hydrostatic balance at the $\epsilon_{\text{mach}}$ level. Qualitatively, figures \Cref{fig:STATIONARY-NORMS} show that both the standard SPH and the potential-temperature formulations are not well-balanced; however, the standard SPH performs considerably worse than the pressure-potential temperature formulation. For standard SPH, the magnitude of the spurious velocity reaches $\approx \qty{e20}{\metre\per\second}$, whereas with the pressure-potential temperature formulation, the magnitudes are not greater than $\approx \qty{e2}{\metre\per\second}$.

\begin{figure}[H]
	\centering
	\begin{subfigure}[b]{\linewidth}
		\centering
		\includegraphics[width=\linewidth]{experiments/stationary_norms_SPH}
		\caption{Classical SPH formulation.}
		\label{fig:STATIONARY-NORMS-SPH}
	\end{subfigure}

	\vspace{0.5cm} 

	\begin{subfigure}[b]{\linewidth}
		\centering
		\includegraphics[width=\linewidth]{experiments/stationary_norms_PTH}
		\caption{Pressure-potential temperature formulation.}
		\label{fig:STATIONARY-NORMS-PTH}
	\end{subfigure}
	\caption{Time evolution of the velocity-field norms $\abs{\vb{v}}_{\infty}$ and $\abs{\vb{v}}_2$ for the stationary hydrostatic balance, comparing each standard formulation with its well-balanced (WB) variant; the vertical axis is logarithmic. For standard SPH \subref{fig:STATIONARY-NORMS-SPH}, the norms sit at approximately \qtyrange{10}{15}{\metre\per\second} until $t \approx \qty{50}{\second}$ and then diverge rapidly; the standard pressure-potential temperature formulation \subref{fig:STATIONARY-NORMS-PTH} stays bounded but nonzero. 
	In both cases, the WB variant maintains the norms at machine epsilon throughout.}
	\label{fig:STATIONARY-NORMS}
\end{figure}

The plots \Cref{fig:STATIONARY-BLOWUPS} show the spurious velocity field at the instant the maximum particle velocity reaches about half the speed of sound (Mach number of $\approx 0.5$), which is roughly $\abs{\vb{v}} \approx \qty{120}{\meter\per\second}$. This threshold is chosen such that we can clearly make a distinction between numerical noise in the simulation and a serious breakdown of the hydrostatic equilibrium.
Notice the velocity field at the top part of the computational domain. During the development of the simulation code, an issue with the top boundary condition has been faced --- if the top of the domain was chosen as open, in the very first instants of the simulation the whole system abruptly collapsed. The reason is that with an open lid, the particles near the boundary have \emph{vacancies in their kernel support} --- they are influenced by an \emph{artificial vacuum}. This issue is common in SPH and the standard practice is to add layers of \emph{ghost particles} \autocite{zhangBoundaryConditionsWeaklycompressible2025,adamiGeneralizedWallBoundary2012}. The result of introducing ghost particles, however, is the \emph{reflection} at the top boundary, which is clearly visible in the flow field.

From \Cref{fig:STATIONARY-STILNESS} it is evident that the perturbation SPH and the well-balanced pressure-potential temperature formulations \emph{preserve the hydrostatic balance exactly.} No spurious velocity field is present.

\begin{figure}[H]
	\centering
	\begin{subfigure}[b]{\linewidth}
		\centering
		\includegraphics[width=\linewidth]{experiments/standard_SPH}
		\caption{Classical SPH formulation.}
		\label{fig:STATIONARY-BLOWUP-SPH}
	\end{subfigure}

	\vspace{0.5cm} 

	\begin{subfigure}[b]{\linewidth}
		\centering
		\includegraphics[width=\linewidth]{experiments/standard_PTH}
		\caption{Standard pressure-potential temperature formulation.}
		\label{fig:STATIONARY-BLOWUP-PTH}
	\end{subfigure}
	\caption{The flow field blowup at $t \approx \qty{50}{\second}$. The vectors show the emergence of non-physical motions caused by the failure to preserve the hydrostatic balance. The displacement of the particles further creates spurious forces, leading to breakdown. Both figures also capture a reflection of the spurious wave at the top boundary.}
	\label{fig:STATIONARY-BLOWUPS}
\end{figure}


\begin{figure}[H]
	\centering
	\begin{subfigure}[b]{\linewidth}
		\centering
		\includegraphics[width=\linewidth]{experiments/well_balanced_SPH}
		\caption{Perturbation SPH formulation.}
		\label{fig:STATIONARY-STILNESS-SPH}
	\end{subfigure}

	\vspace{0.5cm} 

	\begin{subfigure}[b]{\linewidth}
		\centering
		\includegraphics[width=\linewidth]{experiments/well_balanced_PTH}
		\caption{Well-balanced pressure-potential temperature formulation.}
		\label{fig:STATIONARY-STILNESS-PTH}
	\end{subfigure}
	\caption{In both the perturbation/well-balanced formulations, the atmosphere remains static. The hydrostatic balance is exactly preserved.}
	\label{fig:STATIONARY-STILNESS}
\end{figure}

\section{Mountain waves (evolutionary case)}
\label{sec:EVOLUTIONARY}

We have shown in the previous section that both the standard SPH and the pressure-potential temperature formulations fail to preserve hydrostatic balance, whereas the perturbation/well-balanced variants succeed. As we initialize the non-stationary simulations from the hydrostatic balance and wish to capture internal gravity waves, only the \emph{perturbation SPH} and the \emph{well-balanced pressure-potential temperature} formulations have been used in the non-stationary cases.

The panels in \Cref{fig:WAVE-EMERGENCE-PTH} illustrate the emergence and early upward propagation of mountain waves using the well-balanced pressure-potential temperature formulation. As noted in \Cref{sec:FLOW-CORRUGATION}, stationary mountain waves take time to fully develop. Unfortunately, we cannot observe the propagation of these waves further, as the simulation is ultimately disrupted by \emph{tensile instability}, with its onset visible in \subref{fig:pth-wave-60}. The final section is dedicated to discussing this tensile instability in more detail.

\begin{figure}[H]
    \centering
    
    \begin{subfigure}[b]{0.50\linewidth}
        \centering
        \includegraphics[width=\linewidth]{wave_PTH_0035}
        \caption{$t = \qty{420}{\second}$}
        \label{fig:pth-wave-35}
    \end{subfigure}\hfill
    \begin{subfigure}[b]{0.50\linewidth}
        \centering
        \includegraphics[width=\linewidth]{wave_PTH_0040}
        \caption{$t = \qty{480}{\second}$}
        \label{fig:pth-wave-40}
    \end{subfigure}

    \vspace{0.4cm} 

    \begin{subfigure}[b]{0.50\linewidth}
        \centering
        \includegraphics[width=\linewidth]{wave_PTH_0045}
        \caption{$t = \qty{540}{\second}$}
        \label{fig:pth-wave-45}
    \end{subfigure}\hfill
    \begin{subfigure}[b]{0.50\linewidth}
        \centering
        \includegraphics[width=\linewidth]{wave_PTH_0050}
        \caption{$t = \qty{600}{\second}$}
        \label{fig:pth-wave-50}
    \end{subfigure}

    \vspace{0.4cm}

    \begin{subfigure}[b]{0.50\linewidth}
        \centering
        \includegraphics[width=\linewidth]{wave_PTH_0055}
        \caption{$t = \qty{660}{\second}$}
        \label{fig:pth-wave-55}
    \end{subfigure}\hfill
    \begin{subfigure}[b]{0.50\linewidth}
        \centering
        \includegraphics[width=\linewidth]{wave_PTH_0060}
        \caption{$t = \qty{720}{\second}$}
        \label{fig:pth-wave-60}
    \end{subfigure}

    \caption{The emerging mountain wave captured using the well-balanced pressure-potential temperature formulation. The panels illustrate the \emph{upward propagation} of the wave, visualized via the vertical velocity $w$ and lines of constant potential temperature. However, as seen in the final panel \subref{fig:pth-wave-60}, the simulation is ruined due to the onset of \emph{tensile instability}.}
    \label{fig:WAVE-EMERGENCE-PTH}
\end{figure}

As shown in \Cref{fig:WAVE-EMERGENCE-SPH}, in our experiments, the perturbation SPH method fails to capture the mountain wave dynamics. While the well-balanced pressure-potential temperature formulation successfully propagated the wave upward, the perturbation method yields a noisy unstructured flow field. 

We hypothesize that the reason lies in poor conservative properties of the perturbation SPH method. As we mentioned in \Cref{sec:GENERAL-CLASS}, a consistent use of \emph{adaptive smoothing lengths} requires the $\grad h$ terms to be present in the momentum equation. In our case, a simpler approach to $h-$adaptivity has been used together with the perturbation SPH method --- the particle's smoothing length at each time step is updated proportionally to the particle's density. In particular, no $\grad h$ terms have been included in the formulation.\sidenote{It is also not straightforward how the terms could be included. Recall that they appeared naturally when deriving the SPH equations from the discrete Lagrangian in \Cref{sec:GENERAL-CLASS}. However, the standard derivation presented in \Cref{sec:CORNERSTONES} does \emph{not account} for variable $h$ at all.}

\begin{figure}[H]
    \centering
    
    \begin{subfigure}[b]{0.50\linewidth}
        \centering
        \includegraphics[width=\linewidth]{wave_SPH_0030}
        \caption{$t = \qty{360}{\second}$}
        \label{fig:wave-30}
    \end{subfigure}\hfill
    \begin{subfigure}[b]{0.50\linewidth}
        \centering
        \includegraphics[width=\linewidth]{wave_SPH_0035}
        \caption{$t = \qty{420}{\second}$}
        \label{fig:wave-35}
    \end{subfigure}

    \vspace{0.4cm} 

    \begin{subfigure}[b]{0.50\linewidth}
        \centering
        \includegraphics[width=\linewidth]{wave_SPH_0040}
        \caption{$t = \qty{480}{\second}$}
        \label{fig:wave-40}
    \end{subfigure}\hfill
    \begin{subfigure}[b]{0.50\linewidth}
        \centering
        \includegraphics[width=\linewidth]{wave_SPH_0045}
        \caption{$t = \qty{540}{\second}$}
        \label{fig:wave-45}
    \end{subfigure}

    \vspace{0.4cm}

    \begin{subfigure}[b]{0.50\linewidth}
        \centering
        \includegraphics[width=\linewidth]{wave_SPH_0050}
        \caption{$t = \qty{600}{\second}$}
        \label{fig:sph-wave-50}
    \end{subfigure}\hfill
    \begin{subfigure}[b]{0.50\linewidth}
        \centering
        \includegraphics[width=\linewidth]{wave_SPH_0055}
        \caption{$t = \qty{660}{\second}$}
        \label{fig:sph-wave-55}
    \end{subfigure}

    \caption{Time evolution of the flow field computed using the perturbation SPH method. The panels display the vertical velocity component $w$ and lines of constant potential temperature. In contrast to the well-balanced pressure-potential temperature, no emergence or upward propagation of mountain waves is observed.}
    \label{fig:WAVE-EMERGENCE-SPH}
\end{figure}

\section{Tensile instability}
\label{sec:INSTABILITY}

The phenomenon observed in our simulation of mountain waves that led to a complete numerical breakdown is well-known in the SPH community and is called \emph{tensile instability} \autocite{morrisStudyStabilityProperties1996,schuesslerCommentsSmoothedParticle1981,phillipsNumericalMethodThreedimensional1985,sunDeltaALESPHModelImproved2017,priceSmoothedParticleHydrodynamics2012a}. Tensile instability is a \emph{geometric instability} that arises when the interaction between particles causes them to \emph{artificially} attract rather than repel, leading to unphysical particle clumps separated by empty voids \autocite{morrisStudyStabilityProperties1996}. In our simulations, the particles lost their regular spacing, introducing yet another spurious unphysical effect that escalated into a total divergence of the velocity field. 

\begin{figure}[H]
	\centering
	\begin{subfigure}[b]{\linewidth}
		\centering
		\includegraphics[width=\linewidth]{crash_SPH_0062}
		\caption{Perturbation SPH formulation, $t = \qty{750}{\second}$.}
		\label{fig:SPH-CRASH-750}
	\end{subfigure}

	\vspace{0.5cm} 

	\begin{subfigure}[b]{\linewidth}
		\centering
		\includegraphics[width=\linewidth]{crash_SPH_0064}
		\caption{Perturbation SPH formulation, $t = \qty{770}{\second}$. }
		\label{fig:SPH-CRASH-770}
	\end{subfigure}
	\caption{Snapshots displaying the tensile instability in the perturbation SPH formulation. Even though this formulation was unable to capture the mountain wave dynamics, the terrain arguably introduces an irregularity in the flow. This loss of regularity might trigger the tensile instability, as visible in the particle clumping (left). The clumping triggers a sudden, explosion-like divergence in the maximum velocity (right).}
	\label{fig:SPH-CRASH}
\end{figure}

As seen in \Cref{fig:PTH-CRASH}, the onset of the instability in the well-balanced PTH formulation is \emph{fundamentally different} from the standard SPH case. Here, one gets the impression that the irregular configuration of the particles is \emph{triggered} directly by the \emph{emerged mountain wave}. The \emph{physical} density and pressure gradients introduced by the upward-propagating wave distort the particle configuration in the upper atmosphere, eventually creating clumps. Conversely, in the perturbation SPH formulation, as shown in \Cref{fig:SPH-CRASH}, the instability is triggered by the accumulation of spurious background errors.

\begin{figure}[H]
	\centering
	\begin{subfigure}[b]{\linewidth}
		\centering
		\includegraphics[width=\linewidth]{crash_PTH_0057}
		\caption{Well-balanced pressure-potential temperature formulation, $t = \qty{690}{\second}$.}
		\label{fig:PTH-CRASH-690}
	\end{subfigure}

	\vspace{0.5cm} 

	\begin{subfigure}[b]{\linewidth}
		\centering
		\includegraphics[width=\linewidth]{crash_PTH_0060}
		\caption{Well-balanced pressure-potential temperature formulation, $t = \qty{725}{\second}$. }
		\label{fig:PTH-CRASH-725}
	\end{subfigure}
	\caption{Snapshots of the tensile instability in the well-balanced pressure-potential temperature formulation. Unlike the perturbation SPH case, the instability is triggered locally near the mountain by physical perturbations in pressure and density caused by the propagating mountain wave. The oscillatory movements distort the particle distribution and lead to a rapid explosion-like divergence.}
	\label{fig:PTH-CRASH}
\end{figure}

To suppress this instability, a standard remedy is the $\delta$-SPH formulation \autocite{sunDeltaALESPHModelImproved2017}. The core of $\delta$-SPH lies in introducing an \emph{artificial diffusion} term in the evolution equations that acts against the forming of clumps, \ie acts repulsively. This procedure regularizes the density distribution, improves the particle distribution and helps in preventing the breakdown.

However, it is not evident to us how to combine the \emph{well-balanced} property with the $\delta-$SPH framework. As discussed in \Cref{sec:WELL-BALANCED-SPH}, the fundamental issue with the standard SPH formulation is its inability to capture the nonlinear pressure (or potential temperature) profile of the initial hydrostatic state. This leads to inaccurate force estimation, which subsequently causes spurious particle movement. Because $\delta$-SPH introduces its own artificial diffusion that alters the density field, a well-balanced $\delta$-SPH scheme would require constructing a diffusion term that strictly \emph{respects} the nonlinear background profile. 

\end{document}
